\subsection{A coupled perfect kernel with admitted and obstructed pro-five branches}

\begin{statement}{D.27}{Proved}{Exact coupled fibres over normalized finite pro-five shadows}
There are explicit two-generated finite groups
\[
 E_0\xrightarrow{\pi_0}G_0\xrightarrow{\rho}C_5,
 \qquad |\ker\pi_0|=2,
 \tag{D.27.1}
\]
whose respective kernels over $C_5$ are perfect, and in which
$C_5$ acts nontrivially on those kernels by permutation of five
simple factors in each of five rows.  The groups have an
order-five, multiplier-one peripheral automorphism
$R\in\operatorname{Aut}(G_0)$ inducing the identity on $C_5$.
Each of $R,R^2,R^3,R^4$ is obstructed at the central kernel in
\textup{(D.27.1)}.

Here is their exact comparison with a prescribed pro-five
component.  Let $L=\langle l_1,l_2\rangle$ be a finite
$5$-group with an epimorphism
$\eta:L\to C_5=\langle t\rangle$ satisfying
$\eta(l_1)=t$, $\eta(l_2)=1$, and put
$l_3=(l_1l_2)^{-1}$.  Let $\beta\in\operatorname{Aut}(L)$
be specified by the peripheral equalities
\[
 \beta(l_1)=l_1,\qquad
 \beta(l_2)=f^{-1}l_2f,\quad \eta(f)=1,
 \qquad\beta(l_3)=h l_3h^{-1}.
 \tag{D.27.2}
\]
The condition on $f$ holds in particular for the usual
normalized multiplier-one form with $f\in L'$.
Define the actual fibre products
\[
 E_L=E_0\times_{C_5}L,
 \qquad G_L=G_0\times_{C_5}L.
 \tag{D.27.3}
\]
The paired source generators give epimorphisms
$\widehat F\twoheadrightarrow E_L\twoheadrightarrow G_L$,
and the coordinate maps $E_L\to L$ and $G_L\to L$ are their
maximal pro-five quotient maps.  The automorphism
\[
 \alpha_-=(R,\beta)\in\operatorname{Aut}(G_L)
 \tag{D.27.4}
\]
has the prescribed three peripheral conjugacy classes and
has no weak peripheral lift to $E_L$.  On the same quotient
and over the same $\beta$, the automorphism
$\alpha_+=(1,\beta)$ has the explicit generating peripheral
lift $(1,\beta)\in\operatorname{Aut}(E_L)$.

The construction includes an explicit nonidentity $\beta$ of
order five on a finite two-generated $5$-group, given in
\textup{(D.27.15)--(D.27.18)} below.  Thus the perfect kernel,
its action by the quotient, and the prescribed pro-five
automorphism can all be nontrivial.  The admitted branch is
proved for every datum \textup{(D.27.2)}; these coupled
quotients therefore do not exclude such a prescribed
pro-five shadow.
\end{statement}

\begin{proof}
Use $S=A_8$, its explicit double cover $\widetilde S$, and
the five triples $(a,b_j,c_j)$ from D.25, with spins
$s_0=-1$ and $s_1=\cdots=s_4=1$.  Index a copy of $S^{25}$
by $(j,k)\in(\Z/5)^2$, where $j$ labels rows and $k$
labels columns.  Put
\[
 P=S^{\{(j,k)\}},\qquad
 (t\cdot p)_{j,k}=p_{j,k-1},\qquad
 G_0=P\rtimes\langle t\rangle.
 \tag{D.27.5}
\]
In one row write $\delta_k(s)$ for the tuple equal to $s$
in column $k$ and to one elsewhere.  The source generators are
\[
 \begin{split}
 g_1&=((\delta_0(a))_{j=0}^4,t),\\
 g_2&=((\delta_0(b_j))_{j=0}^4,1),\qquad
 g_3=(g_1g_2)^{-1}.
 \end{split}
 \tag{D.27.6}
\]
These are literal elements of the semidirect product with
multiplication $(p,t^r)(q,t^s)=(p(t^r\cdot q),t^{r+s})$.

We prove generation and retain the covering map.  The subgroup
$H=\ker(\widehat F\to C_5)$, with $x\mapsto t,y\mapsto1$,
has the Schreier basis
\[
 x^5,\quad y_k=x^kyx^{-k}\quad(0\leq k\leq4).
 \tag{D.27.7}
\]
Applying \textup{(D.27.6)} gives $x^5\mapsto(a)_{j,k}$ and
$y_k\mapsto(\delta_k(b_j))_j$.  In coordinate $(j,k)$ this
is the epimorphism
\[
 x^5\longmapsto a,\qquad
 y_k\longmapsto b_j,\qquad
 y_r\longmapsto1\quad(r\ne k).
 \tag{D.27.8}
\]
For distinct columns the kernels differ already on one of the
$y_r$.  In the same column and different rows they differ by
the relator witnesses of D.25, after replacing its first and
second generators by $x^5$ and $y_k$.  The simple-group
subdirect argument of D.25 consequently makes the image of
$H$ exactly $P$.  Since $g_1$ maps to $t$, the two source
elements generate all of $G_0$.

Define the central product
\[
 P_0^*=\widetilde S^5/
 \{(\epsilon_0,\ldots,\epsilon_4)\in\{1,-1\}^5:
                        \prod_k\epsilon_k=1\},
 \qquad
 \widetilde P=P_0^*\times(S^5)^4.
 \tag{D.27.9}
\]
The kernel of $P_0^*\to S^5$ is $C_2$, identified with the
product of its five central signs.  Column shift preserves
this defining kernel, so it acts on $\widetilde P$.
Set $E_0=\widetilde P\rtimes C_5$ and let
$\pi_0:E_0\to G_0$ be the induced central extension.
Lift $g_1,g_2$ by replacing their row-zero entries $a,b_0$
with their canonical order-seven lifts in $\widetilde S$;
call the resulting elements $X_1,X_2$, and put
$X_3=(X_1X_2)^{-1}$.
The group $\widetilde P$ is perfect: it is a product of a
quotient of the perfect group $\widetilde S^5$ and copies
of the perfect group $S$.  Consequently $E_0^{\rm ab}=C_5$.
The subgroup generated by $X_1,X_2$ maps onto $G_0$; missing
the central kernel would make it an index-two subgroup of
$E_0$, contrary to this abelianization.  This proves the
source epimorphism to $E_0$.

Let $R$ rotate the row index by
$(Rp)_{j,k}=p_{j+1,k}$ and fix $t$.  It commutes with the
column action and is an automorphism of $G_0$ of order five.
It fixes $g_1$ and preserves the conjugacy classes of $g_2,g_3$.
For $g_2$ this is coordinatewise conjugacy of the $b_j$ in $S$.
For $g_3$ it is the cyclic-product comparison in the wreath
product, whose exact conjugator can be described as follows.
For $v=(v_k,t^\epsilon)$ and $v'=(v'_k,t^\epsilon)$, with
$\epsilon=\pm1$, a base conjugator $(d_k,1)$ must satisfy
\[
 d_kv_kd_{k-\epsilon}^{-1}=v'_k,
 \qquad
 d_{k-\epsilon}=(v'_k)^{-1}d_kv_k.
 \tag{D.27.10}
\]
Choose $d_0$ conjugating the two cycle products.  The displayed
recursion determines all five $d_k$; its last equation is
exactly that cycle-product conjugacy, so it is consistent.
In row $j$, the zeroth coordinate of $g_3^5$ is $c_j$.
Those $c_j$ are conjugate, and \textup{(D.27.10)} supplies
the required conjugator in $P$ for every row rotation.

The peripheral base orders are $(35,7,35)$.  To evaluate their
lifts, temporarily replace $b_0$ in the doubled row by $b_j$
and use canonical odd lifts $\widetilde a,\widetilde b_j$.
The first two lifted peripheral orders are $35,7$.
The fifth power of their product has, in its five columns,
conjugates of $\widetilde a\widetilde b_j$; its
thirty-fifth power is the tuple of five signs $s_j$.
Their product in the central quotient \textup{(D.27.9)}
is $s_j^5=s_j$.  Therefore the natural inverse of that
product has order seventy for $j=0$ and order thirty-five
for $j\ne0$.  The original $X_3$ has order seventy.

In a central extension by $C_2$, an odd-order base element
has precisely one odd-order lift and one even-order lift.
Every permitted conjugate above a fixed base element is
therefore determined by its order.  For $R^j$ with $j\ne0$
the first two required lifts are the canonical odd ones, and
the third is the negative of the natural inverse just
computed, in order to have order seventy.  Their product is
the nonidentity element of $\ker\pi_0$.
Lifting the base conjugators from \textup{(D.27.10)} proves
that all three individually required elements are actual
conjugates.  Thus no remaining conjugacy choice can alter
this product.  This proves the obstruction for all four
nonzero powers of $R$.

We now form \textup{(D.27.3)}.  Equations
\textup{(D.27.2)} give $\eta\beta=\eta$ on the two
generators, hence on $L$, so both displayed product
automorphisms are defined.  The source maps are literally
\[
 x\longmapsto(X_1,l_1),\qquad
 y\longmapsto(X_2,l_2),
 \tag{D.27.11}
\]
and their projections to $G_L$.
For surjectivity, put $H_L=\ker(\widehat F\to L)$.
Then $H_L\triangleleft H$ and $H/H_L\simeq\ker\eta$
is a $5$-group.  The image of $H$ in $\widetilde P$ is
all of $\widetilde P$ by the proved source epimorphism.
Its quotient by the image of $H_L$ is a $5$-group and a
quotient of a perfect group, hence trivial.  Thus $H_L$
maps onto $\widetilde P$.  The image in $E_L$ already
maps onto $L$ and contains its entire kernel, proving
surjectivity.  The proof for $G_L$ is identical with $P$.
The respective kernels $\widetilde P,P$ of the coordinate
maps to $L$ are perfect.  Every homomorphism to a finite
$5$-group kills them, and the coordinate map is itself
onto the $5$-group $L$.  This proves exactly the asserted
maximal pro-five quotient maps.

The admitted branch is the literal automorphism
$(e_0,l)\mapsto(e_0,\beta(l))$ of $E_L$.
It fixes $(X_1,l_1)$.  It conjugates the second source
element by $(1,f^{-1})\in E_L$, since $\eta(f)=1$.
For the third, write $\eta(h)=t^r$.  Since
$\rho\pi_0(X_3)=t^{-1}$, the pair
\[
                  (X_3^{-r},h)\in E_L
 \tag{D.27.12}
\]
is defined and conjugates $(X_3,l_3)$ to
$(X_3,\beta(l_3))$: its first coordinate centralizes
$X_3$.  This supplies the required third conjugator with
the correct fibre-product condition.

For the obstructed branch, choose $d_2,d_3\in P$ with
$d_ig_id_i^{-1}=R(g_i)$, as constructed above.
The second and third base conjugators in $G_L$ are
respectively
\[
 (d_2,f^{-1}),\qquad (d_3g_3^{-r},h).
 \tag{D.27.13}
\]
The first has matching trivial $C_5$ coordinates.  The
second has matching coordinate $t^r$, and multiplying by
$g_3^{-r}$ does not change its conjugation of $g_3$.
They prove that $(R,\beta)$ has precisely the required
peripheral conjugacies.  A weak lift to $E_L$ would, upon
projection to $E_0$, be a weak peripheral lift of $R$ for
the source \textup{(D.27.6)}, which has just been excluded.
The same projection excludes every integer power with
exponent prime to five.

This is an actual semidirect-product coupling.  For example,
conjugation by $(1,l_1)$ sends a nonidentity element supported
in column zero of $P$ to column one.  The typed fibre-product
identification is
\[
 G_0\times_{C_5}L\simeq P\rtimes L,
 \quad ((p,\eta(l)),l)\longmapsto(p,l),
 \tag{D.27.14}
\]
with the specified action through $\eta$; the analogous
map for $E_L$ uses $\widetilde P$.  No direct-product
replacement is made.

We give a concrete nonidentity datum \textup{(D.27.2)}.
Let $\mathcal R=\mathbf F_5\langle X,Y\rangle/
(\text{words of length at least }5)$, let $J$ be its
augmentation ideal, and take the finite group
\[
 L=\langle l_1=\exp X,l_2=\exp Y\rangle\leq1+J.
 \tag{D.27.15}
\]
The finite exponential and logarithm use only denominators
$1,2,3,4$ and $1!,\ldots,4!$, all invertible in
$\mathbf F_5$.  The ambient group has $5^{30}$ elements
and exponent five, since $J^5=0$ and $(1+a)^5=1+a^5=1$.
Its degree-one $X$ coefficient gives the epimorphism
$\eta:L\to C_5$ in the statement.
In the Lie bracket of this associative algebra put
\[
 H_2=[X,Y],\qquad F_3=[X+Y,H_2],\qquad
 H_3=[Y,H_2],\qquad \Delta=[F_3,Y].
 \tag{D.27.16}
\]
The identity $\Delta=[H_3,X+Y]$ follows from Jacobi and
$[H_2,H_2]=0$.  The substitution
\[
             X\longmapsto X,\qquad Y\longmapsto Y+\Delta
 \tag{D.27.17}
\]
is an automorphism of $\mathcal R$ of order five: substitution
in the degree-four polynomial $\Delta$ leaves it unchanged.
The polynomial is nonzero, for example its $XXYY$ coefficient
is one.  With $F=\exp F_3$ and $H=\exp H_3$, direct expansion
gives
\[
 \beta(l_1)=l_1,\qquad
 \beta(l_2)=F l_2F^{-1},\qquad
 \beta((l_1l_2)^{-1})=H(l_1l_2)^{-1}H^{-1}.
 \tag{D.27.18}
\]
The last equality follows also by applying the substitution
to $\log((\exp X\exp Y)^{-1})$: its only changed term is
$-Y$, whose change is $-\Delta=-[H_3,X+Y]$.

We verify membership of the two conjugators in $L$ by actual
words.  Write $c=[l_1,l_2]$, $u=[l_1,c]$, $v=[l_2,c]$,
and $d_1=[l_1,u]$, $d_2=[l_2,u]$, $d_3=[l_2,v]$.
The $d_i$ are central at this truncation.  Multiplication
in $\mathcal R$ gives the exact identities
\[
       F=uvd_1^{-1}d_2^{-1}d_3^{-1},\qquad
       H=v d_2^{-3}d_3^{-1}.
 \tag{D.27.19}
\]
For example the degree-four errors in $\log(uv)$ are the
three $\log d_i$ with coefficients $(1,1,1)$, and those in
$\log v$ have coefficients $(0,3,1)$ over $\mathbf F_5$;
higher terms vanish.  Thus $F,H\in L'$, and
\textup{(D.27.18)} proves that the substitution restricts
to an automorphism of $L$ of order five, with the normalized
conjugator $f=F^{-1}$.  It is nonidentity because
$\beta(l_2)-l_2$ has the nonzero leading term $\Delta$.
For this explicit datum $(R,\beta)$ also has order five,
acts nontrivially on both the perfect kernel and the
maximal pro-five quotient, and all four nonzero powers
retain the central obstruction.

The program \texttt{a\_peripheral\_coupled.py} checks all
twenty-five maps \textup{(D.27.8)}, saves all three hundred
pairwise kernel witnesses, and evaluates the wreath
conjugators \textup{(D.27.10)} literally.  It computes the
central product using the Clifford multiplication of D.25,
obtaining the peripheral orders $(35,7,70)$ in row zero
and $(35,7,35)$ in the other rows, and the allowed products
$1,-1,-1,-1,-1$.  Independently, it verifies every identity
\textup{(D.27.16)--(D.27.19)} by truncated noncommutative
polynomials over $\mathbf F_5$, retaining the literal
conjugator words and their exact coefficients.  Its output
and JSON certificate end with
\texttt{PASS\_COUPLED\_WREATH\_AND\_PRO\_FIVE\_PULLBACK}.
\end{proof}
